\subsection{DeepSeek-R1 and Emergent Reasoning}
\label{sec:deepseek-r1}

A central question in the study of large language model (LLM) reasoning is whether
sophisticated problem-solving behaviors can emerge from reinforcement learning alone,
without reliance on human-annotated reasoning demonstrations. DeepSeek-R1
\cite{deepseek2025r1} provides a compelling affirmative answer. Published by
DeepSeek-AI, the work introduces two models---DeepSeek-R1-Zero and
DeepSeek-R1---that together demonstrate the power of Group Relative Policy
Optimization (GRPO) \cite{shao2024deepseekmath} to elicit emergent chain-of-thought
reasoning, self-reflection, and verification behaviors in large-scale language models.

\subsubsection{DeepSeek-R1-Zero: Pure RL without Supervised Fine-Tuning}

DeepSeek-R1-Zero is trained by applying GRPO directly to the DeepSeek-V3-Base
model \cite{deepseek2024v3}, a 671-billion-parameter Mixture-of-Experts architecture,
without any supervised fine-tuning (SFT) stage. The training setup is deliberately
minimal: the model receives a simple template requiring it to produce a reasoning
process enclosed in \texttt{<think>...</think>} tags, followed by a final answer. No
content-specific biases or human-crafted reasoning demonstrations are provided,
ensuring that any emergent reasoning strategies arise purely from the RL optimization
signal.

The reward design for DeepSeek-R1-Zero is entirely rule-based, consisting of two
components:
\begin{itemize}
    \item \textbf{Accuracy rewards:} For mathematics problems, the final numerical
    answer is extracted and compared against the ground truth. For coding tasks,
    generated solutions are executed against a suite of predefined test cases. These
    rewards provide objective, binary feedback on correctness.
    \item \textbf{Format rewards:} The model is incentivized to place its reasoning
    within the designated \texttt{<think>...</think>} tags and to present the final
    answer in the required format.
\end{itemize}
Crucially, no neural reward model is employed. This deliberate design choice avoids
the well-documented problem of reward hacking, where models learn to exploit biases
in learned reward functions rather than genuinely improving task performance. By
relying exclusively on verifiable, rule-based signals, DeepSeek-R1-Zero ensures that
the RL training signal faithfully reflects actual problem-solving ability.

The results are striking. On the AIME 2024 benchmark, the average pass@1 accuracy of
DeepSeek-R1-Zero increases from an initial 15.6\% to 71.0\% over the course of
training.\footnote{The v1 preprint reports 71.0\%; the revised Nature version (v2)
reports 77.9\% after extended training.} With majority-vote consensus decoding
(cons@64), the score reaches 86.7\%---significantly surpassing the average performance
of human competitors. On MATH-500, DeepSeek-R1-Zero achieves 95.9\% pass@1 accuracy. These gains are achieved with a training configuration of 16
sampled outputs per question, a group size of 512, a learning rate of $3 \times
10^{-6}$, and a KL coefficient of 0.001, over approximately 10{,}400 policy update
steps.

\subsubsection{Emergent Behaviors}

Perhaps the most remarkable aspect of DeepSeek-R1-Zero is the spontaneous emergence
of sophisticated reasoning behaviors that were never explicitly taught. Three
categories of emergent behavior are particularly noteworthy:

\paragraph{Spontaneous chain-of-thought.} As training progresses, DeepSeek-R1-Zero
naturally learns to allocate more ``thinking time'' to problems, with the average
response length increasing steadily throughout training. The model progressively
generates longer reasoning traces---hundreds to thousands of tokens---in which it
explores multiple solution strategies, decomposes complex problems into sub-steps, and
systematically works through intermediate calculations.

\paragraph{Self-reflection and verification.} The model develops the capacity to
recognize errors in its own reasoning and backtrack. It spontaneously produces phrases
such as ``Wait, let me reconsider'' or ``Let me verify this step,'' reassessing
intermediate conclusions and exploring alternative approaches when an initial strategy
appears flawed.

\paragraph{The ``aha moment.''} The authors highlight a particularly striking
phenomenon observed during training: an intermediate checkpoint of DeepSeek-R1-Zero
begins producing the word ``Wait'' with markedly increased frequency during
reflections, signaling a sudden qualitative shift in reasoning patterns. In one
illustrative example, the model pauses mid-derivation and writes: ``\textit{Wait,
wait. Wait. That's an aha moment I can flag here. Let's reevaluate this
step-by-step}.'' The authors describe this as ``an aha moment for us, allowing us to
witness the power and beauty of reinforcement learning''---a moment where the model,
without any explicit instruction, learns to pause, reconsider, and redirect its
reasoning.

These emergent behaviors demonstrate that reinforcement learning with appropriate
incentive structures can give rise to complex cognitive strategies without explicit
supervision, providing evidence that reasoning capabilities can be ``incentivized''
rather than ``taught.''

\subsubsection{Limitations of R1-Zero}

Despite its impressive reasoning capabilities, DeepSeek-R1-Zero exhibits several
practical shortcomings. First, its outputs suffer from \textbf{poor readability}: the
reasoning traces, while effective, are often verbose, disorganized, and difficult for
human readers to follow. Second, the model exhibits pervasive \textbf{language
mixing}, frequently combining English and Chinese within a single chain-of-thought
response---a consequence of DeepSeek-V3-Base being pre-trained on multilingual
corpora. Third, because the RL training focuses exclusively on reasoning tasks with
rule-based rewards, DeepSeek-R1-Zero has \textbf{limited generalization} to broader
capabilities such as creative writing, open-domain question answering, and
instruction following.

\subsubsection{The DeepSeek-R1 Multi-Stage Pipeline}

To address the limitations of R1-Zero while preserving its reasoning strengths,
DeepSeek-R1 is trained through a carefully designed four-stage pipeline:

\paragraph{Stage 1: Cold-start SFT.} Several thousand curated long
chain-of-thought examples are collected to fine-tune the DeepSeek-V3-Base model. These
examples are designed to exhibit a conversational, human-aligned thinking process
presented in first-person perspective, with explicit problem comprehension, detailed
step-by-step reasoning, and integrated reflection and verification. This cold-start
data provides a well-formatted foundation for the subsequent RL stage, producing the
intermediate checkpoint DeepSeek-R1-Dev1.

\paragraph{Stage 2: Reasoning-focused RL.} RL training is applied to the cold-start
checkpoint using GRPO with rule-based accuracy and format rewards, augmented by a
\textbf{language-consistency reward}. This additional reward signal is computed as the
proportion of words in the chain-of-thought that match the target language of the
query, directly penalizing language mixing:
\begin{equation}
    R_{\text{language}} = \frac{\text{Num}(\text{Words}_{\text{target}})}{\text{Num}(\text{Words}_{\text{total}})}.
\end{equation}
The language-consistency reward is added directly to the final reward for both
reasoning and non-reasoning data. Ablation experiments confirm that while this
alignment introduces a slight degradation in raw benchmark performance, it
substantially improves readability and user experience. This stage yields
DeepSeek-R1-Dev2.

\paragraph{Stage 3: Rejection sampling and SFT.} The Dev2 checkpoint is used to
generate reasoning trajectories via rejection sampling, filtering for correctness and
quality. These reasoning samples are combined with non-reasoning data (including
writing, open-domain QA, and instruction-following tasks) to construct a diverse SFT
dataset of approximately 800{,}000 samples. Fine-tuning on this mixed dataset produces
DeepSeek-R1-Dev3, which excels not only in reasoning but also in general language
generation.

\paragraph{Stage 4: RL across all scenarios.} A final RL stage is applied using a
combination of rule-based rewards for reasoning tasks, model-based preference rewards
(helpfulness and safety) for general tasks, and language-consistency rewards. This
comprehensive RL stage uses diverse prompt distributions spanning mathematics, coding,
STEM, logic, and general instruction-following, producing the final DeepSeek-R1 model.
The integration of multiple reward signals enables the model to maintain strong
reasoning while aligning with human preferences for helpfulness and harmlessness.

\subsubsection{Benchmark Results}

Table~\ref{tab:deepseek-r1-results} presents a comparison of DeepSeek-R1-Zero,
DeepSeek-R1, and OpenAI o1-1217 on key reasoning benchmarks. DeepSeek-R1 achieves
performance broadly comparable to OpenAI o1-1217, surpassing it on MATH-500 while
achieving near-parity on AIME 2024. On GPQA Diamond, a benchmark of PhD-level science
questions, DeepSeek-R1 narrows the gap to within 4.2 percentage points. These results
are particularly significant given that DeepSeek-R1 is open-weight and trained at a
fraction of the estimated cost of proprietary alternatives.

\begin{table}[ht]
\centering
\caption{Benchmark comparison of DeepSeek-R1-Zero, DeepSeek-R1, and OpenAI o1-1217.
All figures are pass@1 (\%) except where noted. Results for OpenAI o1-1217 are from
official reports. AIME 2024 and GPQA Diamond use $k{=}64$ samples; MATH-500 uses
$k{=}16$.}
\label{tab:deepseek-r1-results}
\small
\begin{tabular}{@{}p{0.33\textwidth}ccc@{}}
\toprule
\textbf{Benchmark} & \textbf{R1-Zero} & \textbf{R1} & \textbf{o1-1217} \\
\midrule
AIME 2024 (pass@1)       & 77.9  & 79.8  & 79.2 \\
MATH-500 (pass@1)        & 95.9  & 97.3  & 96.4 \\
GPQA Diamond (pass@1)    & 75.8  & 71.5  & 75.7 \\
\midrule
AIME 2024 (cons@64)      & 86.7  & ---   & ---  \\
Codeforces (percentile)  & 80.4  & 96.3  & 96.6 \\
LiveCodeBench (pass@1)   & 50.0  & 65.9  & 63.4 \\
\bottomrule
\end{tabular}
\end{table}

\subsubsection{Distillation to Smaller Models}

A key practical contribution of DeepSeek-R1 is the demonstration that its reasoning
capabilities can be effectively transferred to substantially smaller models through
knowledge distillation. The authors distill DeepSeek-R1's outputs into six student
models ranging from 1.5 billion to 70 billion parameters, built on the Qwen-2.5 and
Llama-3 model families. The distillation procedure is straightforward: approximately
800{,}000 samples of reasoning trajectories generated by DeepSeek-R1 are used to
fine-tune each student's base model for 2--3 epochs, without any additional RL stage.

The distilled models achieve remarkable performance relative to their size. For
instance, DeepSeek-R1-Distill-Qwen-1.5B---a model with only 1.5 billion
parameters---achieves 28.9\% pass@1 on AIME 2024 and 83.9\% on MATH-500, surpassing
GPT-4o (9.3\% and 74.6\%, respectively) and Claude-3.5-Sonnet (16.0\% and 78.3\%).
The larger DeepSeek-R1-Distill-Qwen-32B reaches 72.6\% on AIME 2024 and 94.3\% on
MATH-500.

A particularly important finding is that \textbf{distillation outperforms direct RL on
smaller models}. When Qwen-2.5-32B-Base is trained with large-scale RL from scratch
(over 10{,}000 steps), the resulting model (Qwen2.5-32B-Zero) achieves 47.0\% pass@1
on AIME 2024 and 91.6\% on MATH-500. In contrast, DeepSeek-R1-Distill-Qwen-32B,
which uses only SFT on distilled data, achieves 72.6\% and 94.3\% on the same
benchmarks---a substantial margin. This result suggests that for models below a
certain capacity threshold, distilling the reasoning patterns of a more powerful
teacher is both more effective and more computationally efficient than attempting to
discover those patterns through RL from scratch. However, the authors note that
advancing beyond the frontiers of existing reasoning capability ultimately requires
large-scale RL on sufficiently powerful base models.

\subsubsection{Significance for GRPO Scaling}

DeepSeek-R1 represents the most prominent demonstration to date of GRPO operating at
the scale of hundreds of billions of parameters. Several aspects are directly relevant
to the study of GRPO scaling. First, the success of R1-Zero establishes that GRPO
with purely rule-based rewards is sufficient to elicit complex emergent reasoning in
large models, without requiring neural reward models or supervised demonstrations.
Second, the observation that smaller base models (7B and 16B) failed to exhibit
meaningful reasoning improvements under the same RL protocol highlights a critical
\textit{scale dependence}: the effectiveness of GRPO-based reasoning RL appears
contingent on the base model possessing sufficient capacity to leverage long
chain-of-thought reasoning. Third, the distillation results reveal an asymmetry in
how GRPO scaling interacts with model size---while large models benefit directly from
RL, smaller models benefit more from inheriting the reasoning patterns of RL-trained
teachers. These findings collectively position DeepSeek-R1 as a foundational reference
point for understanding how GRPO-based reinforcement learning scales across model
sizes, training regimes, and task domains.
